\par\clearpage\SourceLeaf{304}\ArithmeticHead{219}{R}{}
\HistoricalFoot{}{28\textsuperscript{*}}
\begin{center}\textbf{Berichtigungen zu den Bänden III und X,1 der Werke.}\end{center}
\begingroup
\fontsize{9.2}{13}\selectfont
\newcommand{\CE}[2]{\par\noindent\hangindent=3.8em\hangafter=1\makebox[3.8em][r]{#1\enspace}#2\par}
\noindent\textbf{Zu Band III:}
\CE{S. 130}{Gl. (4) und (11) fehlt im letzten Gliede der Faktor \(x\), und in (11) hat dieses Glied das Vorzeichen \(-\) statt \(+\).}
\CE{133}{Gl. (23) steht \(\alpha-1\) statt \(\alpha\) und \(\beta-1\) statt \(\beta\).}
\CE{369}{2. Textzeile steht sint statt sive.}
\CE{372}{Zeile 10 v. u. ist simulacque zu streichen.}
\CE{373}{Zeile 9 und 8 v. u. steht \(\displaystyle\frac{x-x'}x\) statt \(\displaystyle\frac{x-x'}{x'}\).}
\CE{377}{Zeile 5 steht \(\displaystyle\log 2\frac{{}^n a}{{}^n b}\) statt \(\displaystyle\log 4\frac{{}^n a}{{}^n b}\).}
\CE{380}{Zeile 1 steht \(\displaystyle\frac1{2^\pi}\) statt \(\displaystyle\frac1{2^n}\).}
\CE{386}{Zeile 10 v. u. steht \(h\) statt \(\displaystyle\frac1h\).}
\CE{405}{Zeile 11 im Nenner steht \(1+2ss+s^4\) statt \(1+2ss-s^4\).}
\CE{420}{Zeile 5 im Nenner von \(A\) steht \(2^{\frac72}\) statt \(2^{\frac74}\).}
\CE{423}{Zeile 8 v. u. im Nenner steht \(\cos\varphi\) statt \(\sin\varphi\).}
\CE{433}{Letzte Formel im Zähler steht \(9e^{-4\frac{\varpi'}\varpi\pi}\) statt \(9e^{-9\frac{\varpi'}\varpi\pi}\).}
\CE{437}{Zeile 11 steht \(-i\sin(2k+2\omega)\omega'\pi\) statt \(+i\sin(2k+2\omega)\omega'\pi\).}
\CE{}{Zeile 13 steht \(-3e\, .\, .\, .\) statt \(-e\, .\, .\, .\)}
\CE{462}{Zeile 10 v. u. steht \(at(1+y^{n-4})\) statt \(at^2(1+y^{n-4})\).}
\CE{478}{Zeile 1 steht \(\displaystyle\frac mn\) statt \(\displaystyle\frac nm\).}
\CE{489}{Zeile 7 v. u. steht \(\displaystyle\frac G{G''}\cos\varphi\) statt \(\displaystyle\frac G{G'}\cos\varphi\).}
\CE{490}{Zeile 7 v. u. steht \(1-hk\sin\varphi^2\) statt \(1-kk\sin\varphi^2\).}
\CE{}{Zeile 5 v. u. 4. Spalte lies proportional.}
\par\medskip\noindent\textbf{Zu Band X,1:}
\CE{S. 16}{muss die Reihe für \(a^\infty\) lauten}
\[
a^\infty=a+\frac23(a'-a)-\frac2{3\mathpunct{.}15}\frac{(a'-a)^2}a
+\frac{14}{9\mathpunct{.}45}\frac{(a'-a)^3}{a^2}- .\, .\, .
\]
\CE{150}{Gleichung (3) steht \(\displaystyle\frac\pi2\) statt \(\displaystyle\frac\pi4\).}
\CE{}{Zeile 6 und 7 sind die Worte klein, gross zu vertauschen.}
\CE{}{Fussnote \textsuperscript{*}) steht 412 statt 413.}
\CE{165}{muss in der Reihe für \(M\) der Koeffizient von \(x^5\) lauten \(-\displaystyle\frac1{60}\).}
\CE{186}{Zeile 7 steht 36384 statt 16384.}
\par\clearpage\SourceLeaf{305}\ArithmeticHead{220}{L}{L. SCHLESINGER, ÜBER GAUSS' FUNKTIONENTHEORETISCHE ARBEITEN.}
\CE{S. 188}{steht die Gleichungsnummer [24] einmal statt [23],}
\CE{}{vorletzte Zeile steht 8024 statt 1024.}
\CE{192}{letzte Zeile steht \((\frac12)^2\) statt \((\frac12)^3\).}
\CE{196}{Gl. [17] steht \(\displaystyle\frac{\sqrt{2\cos v\mathpunct{.}M\cos v}}{{}^{4}\!\surd\sin v}\) statt \(\displaystyle\frac{\sqrt{\cos v\mathpunct{.}M\cos v}}{\surd2\,{}^{4}\!\surd\sin v}\)}
\CE{}{und das zweite Mal \(\sin\varphi\) statt \(\sin5\varphi\).}
\CE{202}{letzte Zeile lies \(R=-z^4Q,\ S=-z^6R\).}
\CE{205}{in dem Ausdruck für \(Qn\varpi\) steht \(\sin n\pi\) statt \(\sin n\pi^2\).}
\CE{}{Zeile 3 v. u. im Nenner steht \(+\displaystyle\frac2{x^8}\) statt \(-\displaystyle\frac2{x^8}\).}
\CE{208}{Zeile 5 in der Klammer fehlt hinter 2 der Faktor \(\displaystyle\frac{x'}{x''}\) und hinter 4 der Faktor \(\displaystyle\frac{x'}{x'''}\).}
\CE{210}{Zeile 5 v. u. steht \(z\) statt 2.}
\CE{211}{Zeile 8 v. u. steht \(B'\) statt \(zB'\).}
\CE{214}{Gl. 4 steht \(\displaystyle\frac{ch\, .\, .\, .}{4x}\) statt \(\displaystyle\frac{4x^2h\, .\, .\, .}c\).}
\CE{223}{2. Spalte, Zeile 3 steht \(e^{-\frac{uu}{4\pi t}}\) statt \(e^{-\frac{\pi uu}t}\).}
\CE{224}{letzte Zeile steht \(M\) statt \(Mi\).}
\CE{225}{Zeile 7 v. u. steht \(e^h\) statt \(e^{-h}\).}
\CE{227}{Zeile 8 steht \(\sin2T'',\ \cos2T''\) statt \(\sin4T'',\ \cos4T''\).}
\CE{253}{Zeile 12 steht \(\displaystyle\lim_{n\to\infty}\frac{b_{-n}}{2^n}\) statt \(\displaystyle\lim_{n\to\infty}\frac{c_{-n}}{2^n}\).}
\CE{254}{Zeile 6 ist linker Hand das \(-\) zu streichen.}
\CE{}{Gl. (19) statt des zweiten \(k^2\) lies \(k^4\).}
\CE{261}{Zeile 11 steht III statt IV.}
\CE{273}{Zeile 18 steht 1798 statt 1799.}
\CE{282}{Zeile 18 v. u. lies \(\gamma it\) statt \(\gamma t\).}
\CE{292}{in der ersten Zeile des Systems [\(2'\)] steht \([xx]^2\) statt \([xx]^3\).}
\CE{293}{Zeile 6 steht \([x^3]\) statt \([x]^3\).}
\CE{298}{Zeile 9 steht \(\displaystyle\left(1-\frac{x^{19}}y\right)\) statt \(\displaystyle\left(1+\frac{x^{19}}y\right)\).}
\CE{301}{Zeile 9 v. u. steht \(+x^{24}\) statt \(-x^{24}\).}
\CE{308}{Zeile 11 v. u. steht \(x^4\) statt \(x^8\).}
\CE{309}{Zeile 2 steht \(\frac92\) statt \(\frac94\).}
\CE{}{Zeile 4 v. u. steht \(ab'+ba'\) statt \(ab'-ba'\).}
\CE{322}{Zeile 2 v. u. steht 415 statt 445.}
\CE{330}{Fussnote letzte Zeile lies: Werke X, 2, Abh. 1, S. 62.}
\CE{394}{Zeile 6 steht kleiner statt grösser.}
\CE{}{Zeile 7 steht kleinsten statt grössten.}
\CE{444}{vorletzte Textzeile und}
\CE{445}{erste Zeile der Fussnote steht 27 statt 24.}
\par\bigskip\begin{center}\rule{.23\linewidth}{.3pt}\end{center}
\par\clearpage\endgroup
